% from ega4-16-fr.tex
\oldpage[IV-4]{52}
dans $B\otimes_A B$) (0, 19.3.5, (iii)); donc $B\otimes_A B$ est aussi une
$A$-algèbre formellement lisse pour les topologies discrètes (0, 19.3.5,
(ii)). Posant $\mathfrak{I}=\mathfrak{I}_{B/A}$, il en résulte que
$B\otimes_A B$ est aussi une $A$-algèbre formellement lisse pour la topologie
$\mathfrak{I}$-préadique (0, 19.3.8); comme par hypothèse
$B=(B\otimes_A B)/\mathfrak{I}$ est une $A$-algèbre formellement lisse pour
les topologies discrètes, la proposition résulte de l'équivalence de a) et b)
dans (0, 19.5.4).

\begin{proposition}[16.10.3]
\label{IV.16.10.3}
Pour qu'un morphisme $f:X\to S$ soit différentiellement lisse, il faut et il
suffit que pour tout $x\in X$, il existe un voisinage ouvert affine $U$ de
$x$, d'anneau $A$, tel que $\Gamma(U,\mathscr{P}_{X/S}^{\infty})$ soit une
$A$-algèbre topologique augmentée isomorphe à l'algèbre complétée $\widehat B$,
où $B=\mathbf{S}_A(V)$, $V$ étant un $A$-module projectif et $B$ étant muni de
la topologie $B^+$-préadique (où $B^+$ est l'idéal d'augmentation). Si
$\Omega_{X/S}^1$ est localement libre de rang fini, on peut remplacer
$\widehat B$ par l'algèbre de séries formelles $A[[T_1,\ldots,T_n]]$.
\end{proposition}

La notion de morphisme différentiellement lisse étant évidemment locale sur
$X$, on peut se borner au cas où $S=\operatorname{Spec}(B)$,
$X=\operatorname{Spec}(C)$. Considérons $C\otimes_B C$ comme une $C$-algèbre
(pour le premier facteur); posons $\mathfrak{I}=\mathfrak{I}_{C/B}$ et
munissons $C\otimes_B C$ de la topologie $\mathfrak{I}$-préadique; on peut
appliquer à la $C$-algèbre topologique $C\otimes_B C$ et à l'idéal
$\mathfrak{I}$ de $C\otimes_B C$ l'équivalence de b) et c) dans (0, 19.5.4),
puisque $(C\otimes_B C)/\mathfrak{I}=C$ est évidemment une $C$-algèbre
formellement lisse pour les topologies discrètes. La topologie sur
$\Gamma(U,\mathscr{P}_{X/S}^{\infty})$ est évidemment la topologie de limite
projective de cet anneau (16.1.11).

On notera que l'entier $n$ de l'énoncé de (16.10.3) est le \emph{rang} de
$\Omega_{X/S}^1$ au point $x$. Nous verrons plus loin (17.13.5) que lorsque
$f$ est différentiellement lisse et localement de type fini, $n$ est aussi
égal à la dimension de la fibre $f^{-1}(f(x))$ au point $x$.

\begin{proposition}[16.10.4]
\label{IV.16.10.4}
Soient $f:X\to S$, $g:S'\to S$ deux morphismes, et posons
$X'=X\times_S S'$, $f'=f_{(S')}:X'\to S'$.
\begin{enumerate}
  \item[(i)] Si $f$ est différentiellement lisse, il en est de même de $f'$.
  \item[(ii)] Inversement, si $g$ est fidèlement plat et quasi-compact, et si
  $f'$ est différentiellement lisse et $\Omega_{X'/S'}^1$ un
  $\mathcal{O}_{X'}$-Module de type fini, $f$ est différentiellement lisse et
  $\Omega_{X/S}^1$ est un $\mathcal{O}_X$-Module de type fini.
\end{enumerate}
\end{proposition}

En effet, si $f$ est différentiellement lisse, les
$\mathscr{G}r_n(\mathscr{P}_{X/S})$ sont des $\mathcal{O}_X$-Modules
\emph{plats}; par suite (16.4.6), l'homomorphisme
\[
  \mathscr{G}r_n(\mathscr{P}_{X/S})\otimes_{\mathcal{O}_X}\mathcal{O}_{X'}
  \longrightarrow\mathscr{G}r_n(\mathscr{P}_{X'/S'})
\]
est bijectif pour tout $n$, et en vertu de la commutativité du diagramme
(16.2.1.3), il résulte de la définition (16.10.1) que $f'$ est
différentiellement lisse. D'autre part, si $g$ est fidèlement plat et
quasi-compact, il résulte encore de (16.4.6) que
\[
  \mathscr{G}r_n(\mathscr{P}_{X/S})\otimes_{\mathcal{O}_X}\mathcal{O}_{X'}
  \longrightarrow\mathscr{G}r_n(\mathscr{P}_{X'/S'})
\]
est bijectif pour tout $n$. Supposons alors $f'$ différentiellement lisse et
$\Omega_{X'/S'}^1$ de rang fini. Comme la projection canonique $X'\to X$ est
un morphisme fidèlement plat et quasi-compact, il résulte d'abord de (2.5.2)
que $\Omega_{X/S}^1$ est un $\mathcal{O}_X$-Module localement libre de rang
fini, puis de (2.2.7) que l'homomorphisme canonique (16.3.1.1) est bijectif,
donc $f$ est différentiellement lisse.

\begin{proposition}[16.10.5]
\label{IV.16.10.5}
Pour qu'un morphisme localement de type fini $f:X\to S$ soit
différentiellement lisse, il faut et il suffit que l'immersion diagonale
$\Delta_f:X\to X\times_S X$ soit quasi-régulière.
\end{proposition}

